%-------------------------------------------------------------------------------
%	SECTION TITLE
%-------------------------------------------------------------------------------
\cvsection{Work Experience}


%-------------------------------------------------------------------------------
%	CONTENT
%-------------------------------------------------------------------------------
\begin{cventries}

%---------------------------------------------------------
  \cventry
    {Senior frontend developer} % Job title
    {\href{https://mercuryo.io/}{Mercuryo}} % Organization with link
    {London, United Kingdom (remote)} % Location
    {Sent. 2025 - Current} % Date(s)
    {
      \begin{cvitems} % Description(s) of tasks/responsibilities
      \vspace{8pt}
      \item[] {\small\textit{Mercuryo is one of the leading companies in crypto-powered payments infrastructure, providing a wallet-integrated on/off-ramp widget, Spend Card, and Wallet that unify fiat and crypto rails for apps and exchanges worldwide (10M users as of May 2025).}}
      \vspace{8pt}
        \item {Owned and maintained Mastercard Crypto Credentials feature frontend, enabling smoother wallet connection, messages/trxs signing, and various user flows. Led the transition of the Crypto credentials to @reown/appkit from @web3-modal/react.}
        \item {Owned end-to-end frontend delivery for Crypto/Spend Card and Ledger Spend integration users (eligibility, KYC handoff, card issuance, top-ups, limits, transaction history, balance, visuals, mocking).}
        \item {Shipped and maintained multiple partner wallet integrations (Ledger, Metamask, TrustWallet and others) of the widget (React+Redux), including white-label embedding, theming, config-driven feature flags, localization, iframes messages communication.}
        \item {Implemented Playwright E2E + vitest unit coverage for high-risk flows (on/off-ramp, KYC redirects, approvals/signatures, card onboarding).}
        \item {Integrated web3 flows into UX: WalletConnect/Reown connection, multiple EVM chain support, message/transaction signing, safe error/retry states across wallets.}
        \item {Improved reliability and debuggability of production issues by integrating client-side observability with Sentry and event analytics.}
        \item {Implemented agents and skills for AI development.}
      \end{cvitems}
    }
    {\vspace{16pt}}

%---------------------------------------------------------
  \cventry
    {Founding frontend engineer} % Job title
    {\href{https://prime.diffuse.fi/}{Diffuse Prime}} % Organization with link
    {Limassol, Cyprus (remote, part-time)} % Location
    {Aug. 2025 - Mar. 2026} % Date(s)
    {
      \begin{cvitems} % Description(s) of tasks/responsibilities
      \vspace{8pt}
      \item[] {\small\textit{Diffuse Prime is a ZK/TEE-powered non-custodial lending protocol for structured leverage deals.}}
      \vspace{8pt}
        \item {Designed architecture and built from scratch nextjs+typescript frontend dApp for the protocol, including wallet connection, on-chain interactions, and user flows for lenders and borrowers. Integrated Sentry, logging and event analytics to monitor the product.}
        \item {In collaboration with the design team, implemented a modern esm UI-kit library with Radix-ui, Tailwind and Tanstack.}
        \item {Implemented testing strategy, including various unit tests with vitest, visual regression tests, lighthouse tests and end-to-end tests with Playwright and local network fixtures. Covered more than 90\% of the codebase with tests.}
        \item {Designed and implemented full cycle CI/CD pipelines for the frontend, ensuring integration and delivery to Vercel.}
        \item {Designed and implemented Nodejs indexer for processing on-chain events and integrated it with nextjs api and Vercel cron jobs.}
        \item {Implemented a custom JavaScript SDK for interacting with the protocol, which was used across the frontend and other tools.}
        \item {Organized monorepo with npm workspaces, implemented automated release and versioning process with changesets, introduced best security practices and audits, and set up various checks in CI.}
        \item {Collaborated closely with founders as a sole frontend engineer to design and implement various features and visual elements, to ensure the product meets the needs of its users.}
        \item {Implemented AI infrastructure for the repository, writing rules for agents and creating agents for various purposes.}
      \end{cvitems}
    }
    {\vspace{16pt}}

%---------------------------------------------------------
  \cventry
    {Frontend developer} % Job title
    {\href{https://nil.foundation}{Nil foundation}} % Organization with link
    {Limassol, Cyprus (remote)} % Location
    {May. 2022 - Aug. 2025} % Date(s)
    {
      \begin{cvitems} % Description(s) of tasks/responsibilities
      \vspace{8pt}
      \item[] {\small\textit{Nil Foundation is an Ethereum L2 startup working on Ethereum horizontal scalability.}}
      \vspace{8pt}
        \item {Built from scratch and released proof market DEX-like interface, which is a decentralized exchange of computation powers for proofs generation.}
        \item {Built a block explorer for the NIL Layer 2 chain with detailed views for blocks, transactions, and accounts. Designed and optimized the frontend architecture for performance and scalability. Also contributed to the Node.js backend powering the explorer.}
        \item {Developed an interactive, sharding-specific smart contract playground for the NIL chain. The tool enables users to write and compile Solidity code, save snippets, deploy contracts, and interact with them — all within a unified interface.}
        \item {Built a custom JavaScript SDK for interacting with the Nil blockchain, tailored to support SSZ serialization and Poseidon hashing. This library became the foundation for multiple critical tools, including a CLI, Hardhat plugin, and browser wallet extension.}
        \item {Developed and maintained the corporate website using Next.js with a Strapi CMS backend, along with additional marketing pages to support product outreach.}
        \item {Created and maintained a reusable React UI component library, adopted across nearly all NIL interfaces to ensure design consistency and development efficiency.}
        \item {Built and maintained CI/CD pipelines using Github actions CI, Ansible and Jinja}
      \end{cvitems}
    }
    {\vspace{16pt}}

%---------------------------------------------------------
  \cventry
    {Frontend developer} % Job title
    {\href{https://infotecs.ru/}{InfoTeCS JSC}}
    {Saint Petersburg, Russia (on-site)} % Location
    {Mar. 2021 - May. 2022} % Date(s)
    {
      \begin{cvitems} % Description(s) of tasks/responsibilities
      \vspace{8pt}
      \item[] {\small\textit{InfoTeCS is a cybersecurity company providing a range of secure communication solutions for corporate clients.}}
      \vspace{8pt}
        \item {Built Remote Work Management System, a big internal React app, where employees could request remote work, track approval pipelines, and submit daily reports. I created interfaces, pages and forms, worked on performance optimizations and managed API workflows with Redux Saga. I also contributed to backend C\# controllers and validations to support frontend logic. }
        \item {Contributed to Employee performance evaluation system - a large-scale web application built with ASP.NET and a legacy frontend stack using jQuery and Backbone.js.}
        \item {Developed React-based Artifact Management App. This internal tool allowed teams to upload, download, and manage build artifacts across multiple systems. All existing artifacts were aggregated in a dynamic responsive table with server-side sorting and filtering, making it easy for engineers to quickly locate, and interact with the files.}
        \item {Led frontend development for a Next.js documentation portal targeting foreign branches, combining static site generation with dynamic internal routing. This portal was a key touchpoint between the company and its foreign customers. One part of the app was public-facing, the other integrated into an internal doc system.}
        \item {Maintained React components library and closely collaborated with Design department to ensure consistent implementation of the design system.}
      \end{cvitems}
    }
    {\vspace{16pt}}

%---------------------------------------------------------
  \cventry
    {Frontend developer} % Job title
    {\href{https://0day.llc/}{0day Technologies}}
    {Moscow, Russia (remote)} % Location
    {Sept. 2020 - Mar. 2021} % Date(s)
    {
      \begin{cvitems} % Description(s) of tasks/responsibilities
      \vspace{8pt}
      \item[] {\small\textit{0day Technologies is a company specializing in the development of automation tools and information security.}}
      \vspace{8pt}
        \item {Designed and developed a Vue2 and Vuetify-based monitoring service, integrating interactive Leaflet maps for real-time data visualization and location tracking.}
        \item {Worked on big Analytics Reactjs application with Python backend, where I created new UI elements, optimized loading time and implemented advanced search across multiple sources.}
      \end{cvitems}
    }
    {\vspace{16pt}}

%---------------------------------------------------------
\end{cventries}
